% Requerimientos del Proyecto Final -- Big Data y Machine Learning con Spark ML
% Transcripción LaTeX del PDF original (docs/requerimientos.pdf).
%
% Compilar con pdfLaTeX:
%   latexmk -pdf requerimientos.tex
%
\documentclass[11pt,a4paper]{article}

\usepackage[T1]{fontenc}
\usepackage[utf8]{inputenc}
\usepackage{tgpagella}          % TeX Gyre Pagella (serif)
\usepackage[spanish,es-noshorthands]{babel}
\usepackage[a4paper,margin=2.5cm]{geometry}
\usepackage{microtype}
\usepackage{xcolor}
\usepackage[most]{tcolorbox}
\usepackage{booktabs}
\usepackage{tabularx}
\usepackage{enumitem}
\usepackage{titlesec}
\usepackage{needspace}
\usepackage[hidelinks]{hyperref}
\hypersetup{pdftitle={Proyecto Final -- Big Data y Machine Learning con Spark ML},pdfsubject={Curso de Big Data}}

% ---------- Tipografía ----------
\renewcommand{\ttdefault}{lmtt}  % monoespaciada Latin Modern

% ---------- Párrafos ----------
\setlength{\parindent}{0pt}
\setlength{\parskip}{0.45em plus 0.15em minus 0.1em}

% ---------- Títulos ----------
\titleformat{\section}{\normalfont\Large\bfseries}{\thesection.}{0.9em}{}
\titlespacing*{\section}{0pt}{1.7ex plus 0.6ex minus 0.2ex}{0.8ex}
\titleformat{\subsection}{\normalfont\large\bfseries}{\thesubsection.}{0.9em}{}
\titlespacing*{\subsection}{0pt}{1.8ex plus 0.5ex minus 0.2ex}{0.65ex}

% ---------- Listas ----------
\newcommand{\sqbullet}{\raisebox{0.25ex}{\rule{0.46ex}{0.46ex}}}
\setlist[itemize]{label=\sqbullet, leftmargin=1.9em, itemsep=0.25em,
                  topsep=0.6em, parsep=0pt, partopsep=0pt}
\setlist[enumerate]{leftmargin=2.1em, itemsep=0.25em, topsep=0.6em,
                    parsep=0pt, partopsep=0pt, label=\arabic*.}

% ---------- Cajas para código y diagramas ----------
\newtcblisting{codebox}{
  listing only, width=\linewidth,
  colback=black!4, colframe=black, boxrule=0.4pt, sharp corners,
  left=6pt, right=6pt, top=5pt, bottom=5pt, boxsep=0pt,
  fontupper=\ttfamily\small, fontlower=\ttfamily\small
}
\newtcblisting{diagram}{
  listing only, width=\linewidth,
  colback=black!4, colframe=black, boxrule=0.4pt, sharp corners,
  left=4pt, right=4pt, top=4pt, bottom=4pt, boxsep=0pt,
  fontupper=\ttfamily\fontsize{7.5}{9}\selectfont,
  fontlower=\ttfamily\fontsize{7.5}{9}\selectfont
}

\begin{document}

% ======================= Portada =======================
\begin{center}
  \vspace*{0.9cm}
  {\LARGE\bfseries Proyecto Final -- Big Data y Machine Learning con Spark ML\par}
  \vspace{1.1em}
  {\large Curso de Big Data\par}
\end{center}

\vspace{0.5em}

% ======================= 1 =======================
\section{Propósito del proyecto}

El proyecto final tiene como objetivo resolver un problema real de
\textbf{detección, clasificación, predicción, regresión o segmentación} utilizando
\textbf{Spark ML} como capa principal de Machine Learning.

El foco estará en diseñar y justificar una solución completa: desde la obtención y
preparación de los datos hasta el entrenamiento, evaluación, interpretación y análisis de
escalabilidad del modelo.

Cada grupo deberá construir una arquitectura donde los datos sean obtenidos mediante
\textbf{una única vía principal de alimentación}, escogida y justificada entre:

\begin{enumerate}[label=\Alph*., leftmargin=2.3em]
  \item \textbf{Spark}: lectura y procesamiento directamente mediante Spark.
  \item \textbf{Dask}: preparación inicial con Dask, persistencia del resultado y posterior
        lectura desde Spark.
  \item \textbf{Kafka}: recepción de eventos mediante Kafka y procesamiento con Spark
        Structured Streaming.
\end{enumerate}

\textbf{Importante:} Spark ML opera sobre DataFrames de Spark. Por tanto, Dask o Kafka
pueden constituir la capa de ingesta o preparación, pero los datos deberán llegar
finalmente a un DataFrame de Spark antes de entrenar o aplicar el modelo.

% ======================= 2 =======================
\section{Continuidad con las tareas anteriores}

El proyecto debe recuperar elementos desarrollados durante el curso y demostrar que el
grupo comprende cómo se conectan las distintas capas de una arquitectura Big Data.

Cuando sea posible, se recomienda continuar con el mismo dominio o conjunto de datos
utilizado en las tareas anteriores. Si el grupo cambia de dominio o dataset, deberá
justificar la decisión.

Cada grupo deberá incluir una breve sección denominada \textbf{``Análisis de datos a
problema de Machine Learning''}, identificando al menos cuatro elementos previamente
trabajados y su transformación en el proyecto.

Ejemplos:

\vspace{0.25em}
\noindent
\small
\renewcommand{\arraystretch}{1.05}
\begin{tabularx}{\textwidth}{@{}>{\raggedright\arraybackslash}p{0.37\textwidth}
                                   >{\raggedright\arraybackslash}X@{}}
  \toprule
  \textbf{Elemento previo} & \textbf{Transformación en el proyecto} \\
  \midrule
  Pregunta analítica & Variable objetivo o decisión que se desea predecir \\
  EDA y distribuciones & Selección y transformación de variables \\
  Outliers o patrones detectados & Señales potenciales para detección o clasificación \\
  Agregaciones o ventanas & Características de entrada para el modelo \\
  Particionamiento de datos & Paralelismo y escalabilidad del entrenamiento \\
  Streaming desde Kafka & Inferencia continua o procesamiento de eventos \\
  \bottomrule
\end{tabularx}
\vspace{0.2em}

% ======================= 3 =======================
\section{Tipos de problemas permitidos}

El problema debe tener una finalidad verificable. Se permiten, entre otros:

\begin{itemize}
  \item clasificación binaria;
  \item clasificación multiclase;
  \item detección de eventos o condiciones;
  \item regresión o predicción de una variable numérica;
  \item predicción basada en variables temporales construidas mediante ventanas o rezagos;
  \item segmentación o clustering, siempre que exista una pregunta concreta que justifique
        la agrupación;
  \item detección de comportamientos atípicos mediante una estrategia justificable
        implementable con las herramientas utilizadas.
\end{itemize}

No se considerará suficiente ejecutar un modelo únicamente porque está disponible en
Spark ML. El grupo deberá justificar por qué el problema requiere Machine Learning, qué
representa la variable objetivo y cuál sería la utilidad de la predicción o detección
obtenida.

% ======================= 4 =======================
\section{Arquitectura}

El informe deberá incluir un diagrama de la arquitectura utilizada.

\needspace{18\baselineskip}
Ejemplos válidos:

\noindent
\begin{minipage}[t]{0.322\textwidth}\vspace{0pt}
\begin{diagram}
Fuente de datos
      |
      v
    Spark
      |
      v
Feature Engineering
      |
      v
   Spark ML
      |
      v
Prediccion / Evaluacion
\end{diagram}
\end{minipage}\hfill
\begin{minipage}[t]{0.322\textwidth}\vspace{0pt}
\begin{diagram}
Fuente de datos
      |
      v
     Dask
      |
      v
Parquet / almacenamiento
      |
      v
    Spark
      |
      v
   Spark ML
\end{diagram}
\end{minipage}\hfill
\begin{minipage}[t]{0.322\textwidth}\vspace{0pt}
\begin{diagram}
Producer
  |
  v
Kafka
  |
  v
Spark Structured Streaming
  |
  v
Transformacion de datos
  |
  v
Modelo Spark ML
  |
  v
Predicciones
\end{diagram}
\end{minipage}

\medskip

Cada grupo deberá escoger \textbf{una} arquitectura principal y justificarla. No se evalúa
positivamente incorporar Dask, Kafka y Spark simultáneamente si no existe una necesidad
real.

% ======================= 5 =======================
\section{Actividades}

\subsection{1. Definición del problema}

Definir claramente:

\begin{itemize}
  \item el problema que se busca resolver;
  \item la unidad de análisis;
  \item la variable objetivo, si corresponde;
  \item qué significa una predicción correcta;
  \item qué decisión o conocimiento aporta el modelo;
  \item las limitaciones principales del problema.
\end{itemize}

La formulación debe estar expresada primero como un problema del dominio y posteriormente
como un problema de Machine Learning.

\subsection{2. Ingesta y preparación}

Implementar una única vía principal de alimentación: Spark, Dask o Kafka.

El grupo deberá describir:

\begin{itemize}
  \item fuente y volumen de datos;
  \item esquema;
  \item variables utilizadas;
  \item tratamiento de nulos;
  \item duplicados;
  \item categorías;
  \item posibles valores anómalos;
  \item transformación necesaria antes del entrenamiento.
\end{itemize}

Si se utiliza Kafka, se deberá explicar además:

\begin{itemize}
  \item qué representa cada evento;
  \item topic utilizado;
  \item key;
  \item número de particiones;
  \item relación entre la tasa de llegada y la capacidad de procesamiento.
\end{itemize}

\subsection{3. Ingeniería de características}

Construir las variables requeridas por el modelo utilizando Spark.

Dependiendo del problema, pueden utilizarse:

\begin{itemize}
  \item variables numéricas;
  \item codificación de variables categóricas;
  \item agregaciones;
  \item variables derivadas;
  \item características obtenidas mediante ventanas;
  \item rezagos temporales;
  \item conteos o tasas;
  \item escalamiento, cuando corresponda.
\end{itemize}

Se debe justificar cada transformación. No se evaluará positivamente incorporar variables sin
explicar su relación con el problema.

\subsection{4. Pipeline de Spark ML}

El procesamiento deberá organizarse mediante los componentes de \texttt{pyspark.ml}.

\needspace{17\baselineskip}
Ejemplo de constructor mínimo:

\begin{codebox}
from pyspark.ml import Pipeline
from pyspark.ml.feature import VectorAssembler

assembler = VectorAssembler(
    inputCols=["x1", "x2", "x3"],
    outputCol="features"
)

pipeline = Pipeline(stages=[
    assembler,
    modelo
])
\end{codebox}

Este ejemplo entrega solamente la estructura. Cada grupo deberá seleccionar sus
transformaciones, variables y modelo de acuerdo con su problema.

\subsection{5. Baseline obligatorio}

Antes de evaluar modelos más complejos, el grupo deberá definir un \textbf{baseline}.

Ejemplos:

\begin{itemize}
  \item predecir siempre la clase mayoritaria;
  \item utilizar una regla simple del dominio;
  \item utilizar la media o mediana en un problema de regresión;
  \item entrenar un modelo lineal o simple como referencia.
\end{itemize}

El objetivo es demostrar que el modelo propuesto mejora una alternativa razonable y no
solamente que obtiene una determinada métrica.

\subsection{6. Modelos}

Evaluar al menos dos alternativas de modelado disponibles en Spark ML cuando el tipo de
problema lo permita.

Ejemplos:

\begin{itemize}
  \item Logistic Regression;
  \item Decision Tree;
  \item Random Forest;
  \item Gradient-Boosted Trees;
  \item Linear Regression;
  \item Generalized Linear Regression;
  \item K-Means.
\end{itemize}

La selección debe estar justificada. No se requiere utilizar todos los algoritmos
disponibles.

\subsection{7. Separación de entrenamiento y evaluación}

Definir una estrategia de separación de datos consistente con el problema.

Se deberá evitar explícitamente el data leakage.

Si los datos poseen una componente temporal relevante, se espera una separación temporal en
lugar de una división aleatoria cuando corresponda.

\subsection{8. Evaluación}

Las métricas deben corresponder al problema.

\textbf{Clasificación:}

\begin{itemize}
  \item precision;
  \item recall;
  \item F1-score;
  \item matriz de confusión;
  \item ROC-AUC y/o PR-AUC cuando corresponda.
\end{itemize}

No se deberá utilizar Accuracy como única métrica cuando exista desbalance de clases.

\textbf{Regresión:}

\begin{itemize}
  \item MAE;
  \item RMSE;
  \item R\textsuperscript{2};
  \item análisis de errores.
\end{itemize}

\textbf{Clustering:}

\begin{itemize}
  \item Silhouette;
  \item caracterización de los grupos;
  \item interpretación de su utilidad.
\end{itemize}

Cada resultado debe seguir la lógica:

\begin{center}
  \texttt{\textbf{pregunta $\rightarrow$ experimento $\rightarrow$ resultado $\rightarrow$ interpretación}}
\end{center}

\subsection{9. Ajuste del modelo (OPCIONAL)}

Realizar al menos una estrategia de ajuste o comparación de hiperparámetros utilizando
herramientas de Spark ML, por ejemplo:

\begin{itemize}
  \item TrainValidationSplit;
  \item CrossValidator;
  \item ParamGridBuilder.
\end{itemize}

El objetivo no es realizar una búsqueda exhaustiva, sino demostrar que el grupo entiende el
efecto de los parámetros sobre el comportamiento del modelo.

\subsection{10. Escalabilidad}

Realizar un experimento que permita observar el comportamiento del pipeline al modificar al
menos una de las siguientes variables:

\begin{itemize}
  \item volumen de datos;
  \item número de particiones;
  \item recursos disponibles;
  \item cantidad de variables;
  \item tasa de eventos, cuando se utilice Kafka.
\end{itemize}

Se deberá registrar el tiempo de ejecución y discutir si el cuello de botella se encuentra
en:

\begin{itemize}
  \item lectura;
  \item preparación;
  \item shuffle;
  \item feature engineering;
  \item entrenamiento;
  \item inferencia;
  \item almacenamiento;
  \item Kafka;
  \item número de particiones;
  \item capacidad del consumidor.
\end{itemize}

\subsection{11. Interpretación del modelo (OPCIONAL)}

El grupo deberá explicar qué aprendió el modelo y qué variables parecen relevantes,
utilizando las herramientas disponibles para el algoritmo seleccionado.

Cuando el algoritmo permita obtener importancia de variables, coeficientes u otra
información interpretable, esta deberá ser analizada.

No basta con reportar una métrica final.

\subsection{12. Si se utiliza Kafka}

Cuando la arquitectura seleccionada utilice Kafka, se deberá distinguir entre:

\begin{itemize}
  \item \textbf{entrenamiento}: realizado sobre un DataFrame de Spark con un conjunto de
        datos definido;
  \item \textbf{inferencia}: aplicación del modelo entrenado sobre los eventos recibidos.
\end{itemize}

Una arquitectura esperada puede ser:

\noindent
\begin{diagram}
Datos historicos
     |
     v
Spark -> Spark ML -> modelo entrenado
                       |
                       v
Kafka -> Structured Streaming -> transformacion -> prediccion
\end{diagram}

\medskip

El grupo deberá explicar cómo se administran los offsets, el checkpointing y qué ocurriría
si el procesamiento se interrumpe temporalmente.

% ======================= 6 =======================
\section{Preguntas mínimas que debe responder el proyecto}

Al finalizar, el grupo deberá ser capaz de responder:

\begin{enumerate}[itemsep=0.55em]
  \item ¿Qué problema real intenta resolver?
  \item ¿Por qué es un problema de Machine Learning?
  \item ¿Cuál es la variable objetivo?
  \item ¿Qué representa cada observación?
  \item ¿Qué variables utiliza el modelo?
  \item ¿Cómo se construyeron esas variables?
  \item ¿Cuál es el baseline?
  \item ¿Qué modelos se compararon?
  \item ¿Qué métrica es la más importante y por qué?
  \item ¿Existe desbalance, leakage o dependencia temporal?
  \item ¿Cuál fue el principal cuello de botella?
  \item ¿Cómo cambia el rendimiento al aumentar la escala?
  \item ¿Qué limitaciones tiene la solución?
  \item ¿Qué decisión podría tomarse utilizando el resultado del modelo?
\end{enumerate}

% ======================= 7 =======================
\section{Entregables}

\begin{enumerate}[itemsep=0.55em]
  \item \textbf{Notebook ejecutable}, con todas las celdas ejecutadas y sus resultados
        visibles. Con problema, metodología, resultados, evaluación, escalabilidad y
        conclusiones.
  \item Diagrama de arquitectura.
  \item Presentación y defensa.
  \item Código o archivos auxiliares necesarios para reproducir la solución.
  \item Declaración de uso de herramientas de Inteligencia Artificial.
\end{enumerate}

% ======================= 8 =======================
\section{Criterios de evaluación}

\noindent
\renewcommand{\arraystretch}{1.15}
\begin{tabularx}{\textwidth}{X}
  \toprule
  \textbf{Criterio} \\
  \midrule
  Formulación y relevancia del problema \\
  Diseño de arquitectura e ingesta \\
  Preparación e ingeniería de características \\
  Uso correcto de Spark ML y Pipeline \\
  Baseline y comparación de modelos \\
  Evaluación y selección de métricas \\
  Análisis de escalabilidad y cuellos de botella \\
  Interpretación y discusión de resultados \\
  Reproducibilidad y calidad de la defensa \\
  \bottomrule
\end{tabularx}

% ======================= 9 =======================
\section{Restricciones}

\begin{itemize}
  \item El modelo principal debe implementarse utilizando \textbf{Spark ML}.
  \item No se acepta reemplazar Spark ML por scikit-learn como núcleo del proyecto.
  \item El proyecto debe utilizar una sola vía principal de ingesta: Spark, Dask o Kafka.
  \item La elección de la arquitectura debe estar justificada por el problema y no por
        acumular tecnologías.
  \item El notebook debe poder ejecutarse de principio a fin.
  \item Las conclusiones deben estar respaldadas por resultados medibles.
  \item Toda herramienta de IA utilizada debe ser declarada junto con la forma en que se
        verificó su resultado.
\end{itemize}

% ======================= 10 =======================
\section{Foco de la defensa}

La defensa no se centrará en recordar sintaxis. El grupo deberá justificar decisiones como:

\begin{itemize}
  \item por qué escogió determinada arquitectura;
  \item por qué una variable puede o no utilizarse como feature;
  \item por qué una métrica es apropiada;
  \item cómo detectó un posible caso de leakage;
  \item por qué un modelo supera al baseline;
  \item qué ocurre al aumentar el volumen de datos;
  \item dónde se encuentra el cuello de botella;
  \item qué modificaría para llevar la solución a un entorno productivo.
\end{itemize}

\end{document}
